\section{Related Literature}
\label{sec:literature}

% Paragraph 1: Career dynamics and path dependence
This paper connects to the literature on career dynamics and path dependence. \citet{Oyer2006_initial_placement} shows that initial labor market conditions have persistent effects on the careers of economists, with initial placement determining long-run research productivity and job quality. \citet{Oreopoulos2012_graduating_recession} finds that graduating during a recession reduces earnings for up to a decade. \citet{Oyer2008_investment_banker} demonstrates that stock market conditions at the time of MBA graduation determine whether graduates enter investment banking, with lasting income consequences. \citet{Chetty2014_where_is_land} documents geographic variation in intergenerational mobility, establishing that the environment in which a person starts shapes their long-run trajectory. The common thread is that initial conditions matter because they determine access to subsequent opportunities---a form of cumulative advantage \citep{Merton1968_matthew_effect}. The present paper tests whether this logic extends to within-career opportunity shocks that are randomly assigned, and separates the access and ability dimensions of the response.

% Paragraph 2: Returns to early opportunities and access shocks
A second strand examines the returns to marginal opportunities and access shocks. The closest labor market analog to the lucky loser setting is \citet{Pallais2014_inefficient_hiring}, who shows that random provision of a single work opportunity on oDesk leads to persistent gains in subsequent employment---a finding that mirrors the access-begets-access mechanism documented here. \citet{Abel2020_reference_letters} find that randomly assigned reference letters improve labor market outcomes for disadvantaged job seekers in South Africa. \citet{Startz2021_interviewer_effects} demonstrates that random interviewer assignment affects employment outcomes, establishing that small procedural shocks at the point of access can have lasting consequences. In science, \citet{Azoulay2014_matthew_effect} document a Matthew effect in which early publication success cascades into future productivity gains, and \citet{Wang2019_early_career_setback} find that scientists who narrowly miss early-career grant funding subsequently outperform those who narrowly succeed. \citet{Azoulay2010_superstar_extinction} show that the death of a superstar scientist reduces the productivity of their collaborators, documenting how access to high-quality peers shapes output. College access at the admission threshold produces long-run earnings gains \citep{Zimmerman2014_college_marginal, Hoekstra2009_flagship_earnings, Goodman2017_college_access}. In sports, \citet{Berger2011_can_losing} show that NBA teams trailing at halftime are more likely to win, providing evidence on motivational responses to adversity. This paper contributes to this growing literature on access shocks by providing literal randomization and the ability to separately measure access and ability responses---a distinction that existing studies of opportunity provision cannot typically make.

% Paragraph 3: Sports quasi-experimental literature
A third strand uses professional sports as laboratories for testing economic theories. \citet{Gauriot2019_fooled_randomness} use close tennis matches to show that the ranking system overrewards luck relative to skill. \citet{GonzalezDiaz2012_performing_best} study whether tennis players perform better in high-stakes moments. \citet{Sunde2009_heterogeneity_tennis} tests tournament theory predictions about heterogeneity and incentives. \citet{Ehrenberg1990_tournaments_golf} examines whether golf tournament prize structures have incentive effects. Most directly related, \citet{Maity2025_early_career_sports} study lucky losers across professional sports and find that early-career setbacks are associated with stronger subsequent performance. Their analysis is correlational: they compare lucky losers to other qualifying losers without accounting for the selection that determines who receives a lucky loser slot. The present paper provides the first causal estimates by exploiting the Grand Slam lottery mechanism.

% Gap statement
The gap this paper fills is the absence of clean causal identification combined with dynamic career tracking. Existing studies of initial conditions use quasi-experimental variation that is typically one-shot (graduating in a recession, narrowly winning an election) and measure outcomes at a single long-run horizon. The lucky loser setting provides literal randomization, dynamic outcomes at multiple horizons from 4 to 52 weeks, and the ability to separately measure access and ability. The finding that access effects persist while ability effects are null on the ATP tour---and marginally positive on the WTA tour---is a distinction that the existing literature cannot make because it typically observes only one dimension of career capital.
